\documentclass{../template/note}

\author{Oliver West}
\title{Probability}
\date{}

\begin{document}

    \maketitle

    \tableofcontents

    \lecture{Lecture 1 -- 23/01/26}

    \section{Fundamentals}

    \begin{definition}
        Let $\Omega$ be a set and $\mathcal F$ is a collection of subsets of $\Omega$. We call $\mathcal F$ a \emph{$\sigma$-algebra} if
        \begin{itemize}
            \item $\Omega \in \mathcal F$;
            \item if $A \in \mathcal F$, then $A^\comp \in \mathcal F$;
            \item if $(A_n)_{n \in \NN}$ is a countable collection of sets in $\mathcal F$, then
            \begin{align*}
                \bigcup A_n \in \mathcal F
            \end{align*}
        \end{itemize}
    \end{definition}
    
    \begin{definition}
        Let $\mathcal F$ be a $\sigma$-algebra on $\Omega$. A function
        \begin{align*}
            \PP: \mathcal F \to [0, 1]
        \end{align*}
        is called a \emph{probability measure} if
        \begin{itemize}
            \item $\PP(\Omega) = 1$;
            \item for any countable disjoint collection $(A_n)_{n \in \NN}$ in $\mathcal F$,
            \begin{align*}
                \PP \left(\bigcup A_n\right) = \sum \PP(A_n)
            \end{align*}
        \end{itemize}
        We call $\PP(A)$ the \emph{probability} of $A$.
    \end{definition}
    
    Then $(\Omega, \mathcal F, \PP)$ is a \emph{probability space}.
    
    \begin{definition}
        The elements of $\Omega$ are called \emph{outcomes}, and the elements of the $\sigma$-algebra $\mathcal F$ are called \emph{events}.
    \end{definition}
    
    Importantly, we talk about probabilities of \emph{events}, never probabilities of outcomes -- the domain of $\PP$ is $\mathcal F$, not $\Omega$.

    \begin{remark*}
        When $\Omega$ is countable, we take $\mathcal F$ to be all subsets of $\Omega$. However, this framework will become useful when considering uncountable $\Omega$.
    \end{remark*}
    
    Let us note some important properties of probability measures that follow from the definition.

    \begin{itemize}
        \item $\PP(A^\comp) = 1 - \PP(A)$;
        \item $\PP(\emptyset) = 0$;
        \item $\PP(A) \leq \PP(B)$ if $A \subseteq B$;
        \item $\PP(A \cup B) = \PP(A) + \PP(B) - \PP(A \cap B)$.
    \end{itemize}
    
    \begin{example}
        \begin{itemize}
            \item \emph{Rolling a fair die}. Then $\Omega = \{1, 2, 3, 4, 5, 6\}$, $\mathcal F = \mathcal P(\Omega)$, and
            \begin{align*}
                \PP(A) = \frac{|A|}{6}
            \end{align*}
            since all outcomes are equally likely.
    
            \item $\Omega = \{\omega_1, \dots, \omega_n\}$, $\mathcal F = \mathcal P(\Omega)$, $\PP(A) = \frac{|A|}{|\Omega|}$ models the experiment of picking an element of $\Omega$ uniformly at random.
            
            \item \emph{Picking balls from a bag}. Suppose we have $n$ balls labelled $\{1, \dots, n\}$, and we pick $k \leq n$ balls uniformly at random without replacement. Then $\Omega = \{A:  A \subseteq \{1, \dots, n\}, |A| = k\}$, so $|\Omega| = \binom n k$ and $\PP(\{\omega\}) = \binom n k ^{-1}$ for any $\omega \in \Omega$.
            
            \item \emph{A deck of 52 cards}. Take a well-shuffled deck of 52 cards. Then $\Omega = \{\text{all permutations of 52 cards}\}$, so $|\Omega| = 52!$.
            
            What is, for example, $\PP(\text{top 2 cards are aces})$? It is $\frac{4 \cdot 3 \cdot 50!}{52!}$.

            \item \emph{Largest digit}. Consider a string of $n$ random digits from 0 up to 9. Then $\Omega = \{0, \dots, 9\}^n$, $|\Omega| = 10^n$.
            
            Consider the events $A_k = \{\text{no digit exceeds $k$}\}$, $B_k = \{\text{largest digit is $k$}\}$. Then
            \begin{align*}
                \PP(A_k) &= \frac{(k+1)^n}{10^n} \\
                \implies \PP(B_k) = \PP(A_k \setminus A_{k-1}) &= \frac{(k+1)^n-k^n}{10^n}
            \end{align*}

            \item \emph{Birthday problem}. There are $n$ people. What is the probability that at least 2 of them share the same birthday?
            
            Assume each birthday is equally likely to be one of $\{1, \dots, 365\}$. Then $\Omega = \{1, \dots, 365\}^n$, $|\Omega| = 365^n$, and $\PP(\{\omega\}) = \frac{1}{365^n}$.

            Let $A = \{\text{at least 2 people have the same birthday}\}$. Instead we count $A^\comp = \{\text{all birthdays different}\}$.

            \begin{align*}
                \PP(A^\comp) &= \frac{365 \cdot 364 \cdots (365-n+1)}{365^n} \\
                \implies \PP(A) &= 1 - \frac{365 \cdot 364 \cdots (365-n+1)}{365^n}
            \end{align*}
            
            Perhaps surprisingly, if $n=22$, $\PP(A) = 0.476$, while if $n=23$, $\PP(A) = 0.507$.
        \end{itemize}
    \end{example}
    
    \subsection{Combinatorial analysis}

    Let us introduce some notions.

    \begin{enumerate}
        \item Let $\Omega$ be a finite set with $|\Omega| = n$. We want to partition $\Omega$ into $k$ disjoint subsets $\Omega_1, \dots, \Omega_k$ with $|\Omega_i| = n_i$ and $\sum n_i = n$. How many ways are there to do this? It is

        \begin{align*}
            M &= \binom n {n_1} \cdot \binom {n-n_1} {n_2} \dots \binom{n-(n_1 + \dots + n_{k-1})}{n_k} \\
            &= \frac{n!}{n_1!n_2! \cdots n_k!} \\
            &= \binom{n}{n_1, n_2, \dots, n_k}
        \end{align*}
        
        This is the \emph{multinomial coefficient}.

        \item Suppose $f: \{1, \dots, k\} \to \{1, \dots, n\}$ (with $k \leq n$). $f$ is \emph{strictly increasing} if $x < y \implies f(x) < f(y)$, and $f$ is \emph{increasing} if $x < y \implies f(x) \leq f(y)$.
        
        How many strictly increasing $f: \{1, \dots, k\} \to \{1, \dots, n\}$ are there? Such a function is determined uniquely by its range, so there are $\binom n k$ of them.
    \end{enumerate}


    
    

    
    
    
    
    
\end{document}
